\chapter{注意力机制优化}
在上一章中，我们学习了 Transformer 的核心机制——自注意力（Self-Attention）。
标准的多头注意力（Multi-Head Attention，MHA）在训练和推理中展现了强大的建模能力，但随着模型规模的增长和序列长度的增加，它也面临着一些实际的效率和资源问题。
本章将介绍几种重要的注意力机制优化方法，它们分别从不同的角度解决这些问题。

\section{GQA}
在第6章中，我们学习了多头注意力机制。在标准的多头注意力中，每个注意力头都有独立的查询矩阵 $\boldsymbol{W}^Q$、键矩阵 $\boldsymbol{W}^K$ 和值矩阵 $\boldsymbol{W}^V$。
假设有 $h$ 个注意力头，每个头的维度为 $d_h = d / h$，那么总的参数量为 $h \times 3 \times d \times d_h = 3d^2$。
在推理阶段，每个头的键和值都需要被缓存下来，以便在生成下一个 token 时复用，这就是所谓的 KV Cache。
KV Cache 的大小与头数 $h$ 和序列长度 $n$ 成正比，当模型规模增大或序列变长时，KV Cache 会占用大量显存。

为了减少 KV Cache 的开销，一种直观的想法是让所有注意力头共享同一组键和值，这就是多查询注意力（Multi-Query Attention，MQA）。
在 MQA 中，只有查询矩阵 $\boldsymbol{W}^Q$ 是多头的，而键矩阵 $\boldsymbol{W}^K$ 和值矩阵 $\boldsymbol{W}^V$ 只有一个头，所有查询头共享这一组键值。
这样，KV Cache 的大小从 $h \times n \times d_h$ 减少到 $1 \times n \times d_h$，减少了 $h$ 倍。
然而，MQA 的表达能力下降明显，因为所有头只能依赖同一组键值来计算注意力，模型的多样性受到了限制。

分组查询注意力（Grouped Query Attention，GQA）是 MQA 和 MHA 之间的折中方案。
在 GQA 中，将 $h$ 个查询头分成 $g$ 组（$1 \leq g \leq h$），每组共享一组键和值。
当 $g = 1$ 时，GQA 退化为 MQA；当 $g = h$ 时，GQA 退化为 MHA。
通常取 $g$ 为 8 或 16，这样 KV Cache 的大小减少到 $g \times n \times d_h$，同时保留了较好的表达能力。

具体来说，假设查询头数为 $h$，键值头数为 $g$，每个查询头 $i$ 属于第 $\lfloor i \cdot g / h \rfloor$ 组，使用该组的键和值来计算注意力。
GQA 被广泛应用于 LLaMA 2、Qwen 等大语言模型中，它在保持模型质量的同时，显著降低了推理时的显存占用。

\section{FlashAttention}
在前面的注意力计算中，我们首先计算查询和键的点积得到注意力分数矩阵 $\boldsymbol{S} = \boldsymbol{Q} \boldsymbol{K}^\mathrm{T} / \sqrt{d_k}$，
然后对 $\boldsymbol{S}$ 应用 softmax 得到注意力权重矩阵 $\boldsymbol{P}$，最后乘以值矩阵 $\boldsymbol{V}$ 得到输出 $\boldsymbol{O} = \boldsymbol{P} \boldsymbol{V}$。
这个过程需要将 $\boldsymbol{S}$ 和 $\boldsymbol{P}$ 这两个 $n \times n$ 的矩阵存储在 GPU 的高带宽内存（HBM）中。

然而，GPU 的内存层次结构存在显著差异：HBM 容量大（如 40GB）但带宽有限（如 1.5 TB/s），
而 SRAM（片上内存）容量小（如 20MB）但带宽极高（如 19 TB/s）。
标准注意力算法的瓶颈不在于计算量，而在于 HBM 的读写次数（IO 复杂度）。
对于长度为 $n$ 的序列，$\boldsymbol{S}$ 和 $\boldsymbol{P}$ 的大小为 $O(n^2)$，将它们从 HBM 读写的开销成为性能瓶颈。

FlashAttention 的核心思想是：通过分块计算（tiling），避免将 $O(n^2)$ 的中间矩阵写入 HBM。
具体来说，将 $\boldsymbol{Q}$、$\boldsymbol{K}$、$\boldsymbol{V}$ 分成若干块（block），每次只加载一小块到 SRAM 中计算，
计算完成后直接将结果写回 HBM，而不需要在 HBM 中存储完整的 $\boldsymbol{S}$ 和 $\boldsymbol{P}$。

这里有一个关键问题：softmax 操作需要知道整行的最大值（用于数值稳定性），而分块计算时只能看到一部分。
FlashAttention 通过维护一个 running maximum 来解决这个问题：在计算每一块时，更新当前已见的最大值，
并对之前的部分进行重新缩放（rescaling）。这样，最终的结果与标准注意力在数学上是精确等价的，没有任何近似。

此外，在反向传播时，FlashAttention 采用重计算（recomputation）策略：前向传播时不存储 $\boldsymbol{S}$ 和 $\boldsymbol{P}$，
反向传播时重新计算它们。虽然增加了计算量，但大幅减少了显存占用（从 $O(n^2)$ 降到 $O(n)$），使得训练更长序列成为可能。
FlashAttention 已成为大模型训练和推理的标准组件，被集成到 PyTorch、DeepSpeed 等主流框架中。

\section{PagedAttention}
在大语言模型的推理服务中，通常需要同时处理多个请求，每个请求的序列长度不同且动态增长。
为了支持自回归生成，每个请求都需要维护自己的 KV Cache。传统的做法是为每个请求预分配一块连续的显存，
大小按照最大可能的序列长度来设定。然而，由于实际序列长度远小于最大值，且难以预测，
这种方式会导致大量的显存浪费和碎片化，据统计，内部碎片和外部碎片可浪费 60\%--80\% 的显存。

PagedAttention 借鉴了操作系统中虚拟内存的分页思想，将 KV Cache 划分为固定大小的块（block），
每个块包含固定数量 token 的键值缓存（如 16 个 token）。这些块不需要在显存中连续存放，
而是通过一个块表（block table）来记录逻辑块到物理块的映射关系。

当生成新的 token 时，系统按需分配新的块，而不是一次性分配大块连续显存。
这样，显存的利用率可以接近 100\%，几乎没有浪费。此外，由于不同请求的 KV Cache 可以共享相同的块
（例如在 beam search 或并行采样中，多个候选序列共享前缀），PagedAttention 还可以通过块拷贝（copy-on-write）
进一步减少显存占用。

PagedAttention 是 vLLM 系统的核心技术，它使得在相同显存下可以处理更多的并发请求，显著提升了推理服务的吞吐量。

\section{线性注意力}
标准注意力的计算复杂度是 $O(n^2 d)$，其中 $n$ 是序列长度，$d$ 是头维度。
这是因为需要计算 $n \times n$ 的注意力分数矩阵 $\boldsymbol{S} = \boldsymbol{Q} \boldsymbol{K}^\mathrm{T} / \sqrt{d_k}$。
当序列长度较短时（如几千），这个复杂度是可以接受的；但当序列长度达到数万甚至数十万时，
$O(n^2)$ 的计算和显存开销 becomes 不可承受。

线性注意力（Linear Attention）通过一种巧妙的变换，将注意力的复杂度从 $O(n^2 d)$ 降低到 $O(n d^2)$。
其核心思想是用一个核函数 $\phi$ 来近似 softmax：
$$
	\text{Attention}(\boldsymbol{Q}, \boldsymbol{K}, \boldsymbol{V}) \approx \phi(\boldsymbol{Q}) \left( \phi(\boldsymbol{K})^\mathrm{T} \boldsymbol{V} \right)
$$
其中 $\phi(\boldsymbol{Q})$ 和 $\phi(\boldsymbol{K})$ 的维度都是 $n \times d$。
关键的变化在于计算顺序：标准注意力是先计算 $\boldsymbol{Q} \boldsymbol{K}^\mathrm{T}$（$n \times n$），再乘以 $\boldsymbol{V}$（$n \times d$）；
而线性注意力利用矩阵乘法的结合律，先计算 $\phi(\boldsymbol{K})^\mathrm{T} \boldsymbol{V}$（$d \times d$），再左乘 $\phi(\boldsymbol{Q})$（$n \times d$）。
这样，避免了 $n \times n$ 矩阵的计算和存储。

从形式上看，线性注意力可以看作是一种 RNN：在序列的每一步 $t$，维护一个 $d \times d$ 的状态矩阵 $\boldsymbol{S}_t = \boldsymbol{S}_{t-1} + \phi(\boldsymbol{k}_t) \boldsymbol{v}_t^\mathrm{T}$，
输出 $\boldsymbol{o}_t = \phi(\boldsymbol{q}_t) \boldsymbol{S}_t$。这与第5章介绍的 RNN 结构非常相似，只是状态更新方式不同。

线性注意力的主要挑战在于如何选择合适的核函数 $\phi$ 来近似 softmax，使得近似误差尽可能小。
常见的方法包括使用泰勒展开、随机特征近似（如 Performer）、或者通过训练学习核函数。
尽管线性注意力在长序列处理上有理论优势，但在实际应用中，由于近似误差和硬件优化等因素，
它通常作为标准注意力的补充，用于特定的长序列场景。
